	\chapter{Projective Spaces}

In this chapter we define projective spaces in two different ways. One is based on gluing affine spaces together and the other is based from the $\Proj$ construction. The two definitions are equivalent.

\section{Projective Spaces as Gluing of Affine Schemes}

In this section we construct the projective space $\PP^n_R$ over a ring $R$ via gluing $n + 1$ copies of the affine space. Let $X_0, X_1, \dotsc, X_n$ be variables. Define, for each $i = 0, \dotsc, n$,
\begin{equation*}
	U_i = \Spec R\left[\frac{X_0}{X_i}, \dotsc, \widehat{\frac{X_i}{X_i}}, \dotsc, \frac{X_n}{X_i}\right],
\end{equation*}
where $\widehat{T}$ means that $T$ is omitted. Therefore $U_i$ is isomorphic to the affine space $\AAA^n_R$. We like to think that everything is happening inside the ring of Laurent polynomials $S = R[X_0, \dotsc, X_n, X_0^{-1}, \dotsc, X_n^{-1}]$ which clearly contains these defining rings of $U_i$ as subrings. The gluing datum for $\PP^n_R$ is as follows.
\begin{enumerate}
	\item The index set is $\Set{0, 1, \dotsc, n}$ and the schemes are the $U_i$.
	\item For each $0 \leq i, j \leq n$ and $i \neq j$, set
	\begin{equation*}
		U_{ij} = D_{U_i}\left(\frac{X_j}{X_i}\right) \subseteq U_i.
	\end{equation*}
	Otherwise set $U_{ii} = U_i$.
	
	\item Define $\varphi_{ii} = \id_{U_i}$; and if $i \neq j$, define
	\begin{equation*}
		\varphi_{ji} \colon U_{ij} \to U_{ji}
	\end{equation*}
	to be the isomorphism corresponding to the identity map
	\begin{equation*}
		R\left[\frac{X_0}{X_i}, \dotsc, \widehat{\frac{X_i}{X_i}}, \dotsc, \frac{X_n}{X_i}\right]_{X_j/X_i} \leftarrow R\left[\frac{X_0}{X_j}, \dotsc, \widehat{\frac{X_j}{X_j}}, \dotsc, \frac{X_n}{X_j}\right]_{X_i/X_j}.
	\end{equation*}
\end{enumerate}
Since the morphisms $\varphi_{ji}$ are defined by equalities the cocycle condition holds trivially for them. The resulting scheme after gluing is called the \inlinedef{projective space of relative dimension $n$ over $R$} and is denoted by $\PP^n_R$. The scheme $U_i$ is regarded as a subscheme of $\PP^n_R$ and is denoted by $D_+(X_i)$.

\begin{proposition}
	The canonical ring homomorphism $R \to \Gamma(\PP^n_R, \mathcal{O}_{\PP^n_R})$ is an isomorphism.
\end{proposition}
\begin{proof}
	If $n = 0$, then $\PP^n_R = \Spec R$ and there is nothing to prove. So let $n \geq 1$. For each $i$, write
	\begin{equation*}
		R_i = R\left[\frac{X_0}{X_i}, \dotsc, \widehat{\frac{X_i}{X_i}}, \dotsc, \frac{X_n}{X_i}\right].
	\end{equation*}
	Firstly, observe that the restriction morphism $\Gamma(U_i, \mathcal{O}_{\PP^n_R}) \to \Gamma(U_{ij}, \mathcal{O}_{\PP^n_R})$ is injective because it is the localisation map
	\begin{equation*}
		R_i = R\left[\frac{X_0}{X_i}, \dotsc, \widehat{\frac{X_i}{X_i}}, \dotsc, \frac{X_n}{X_i}\right] \to R\left[\frac{X_0}{X_i}, \dotsc, \widehat{\frac{X_i}{X_i}}, \dotsc, \frac{X_n}{X_i}\right]_{X_j/X_i}.
	\end{equation*} 
	The element $X_j/X_i$ is a non-zerodivisor of the ring $R_i$. This gives the above observation. Hence we may regard this restriction map as an inclusion of subrings. The sheaf condition with respect to the open cover $\Set{U_i}$ of $\PP^n_R$ implies that
	\begin{equation*}
		\Gamma(\PP^n_R, \mathcal{O}_{\PP^n_R}) = \Set*{(f_i) \in \prod_{i = 1}^n \Gamma(U_i, \mathcal{O}_{\PP^n_R}) : f_i|_{U_{ij}} = f_j|_{U_{ji}} \text{ for all } i, j}.
	\end{equation*}
	By the above observation, $f_i|_{U_{ij}} = f_i$ for all $i, j$. Therefore an element $f \in S$ defines a global section of $\PP^n_R$ if and only if $f \in R_i$ for all $i$. In other words,
	\begin{equation*}
		\Gamma(\PP^n_R, \mathcal{O}_{\PP^n_R}) \cong \bigcap_{i = 1}^n R_i. 
	\end{equation*}
	An element $f \in S$ belongs to $R_i$ if and only if the total degree is zero and each monomial $r X_0^{\alpha_0} \dotsb X_n^{\alpha_n}$ of $f$ with $r \in R$, $\alpha_n \in \ZZ$ satisfies $\alpha_k \geq 0$ for all $k \neq i$. From this it follows that $\Gamma(\PP^n_R, \mathcal{O}_{\PP^n_R}) \cong \bigcap_{i = 1}^n R_i = R$.
\end{proof}

\subsection*{Morphisms into Projective Spaces}
\begin{proposition}
	Let 
\end{proposition}

\begin{corollary}
	If $n > 0$, then $\PP^n_R$ is not affine.
\end{corollary} 
\begin{proof}
	Suppose not. Then $\PP^n_R \cong \Spec R$ by the anti-equivalence of the category of affine schemes and the category of rings (\cref{chap3:antiequivalence-of-affine-schemes-and-rings}).
\end{proof}

\section{Vanishing Schemes}





\section{Basics of Graded Rings}
\subsection*{Graded Rings, Graded Algebras, and Graded Modules}

\begin{definition}~
	\begin{enumerate}
		\item A \inlinedef{graded ring} is a ring with a decomposition $A = \oplus_{d \geq 0} A_d$ of abelian groups such that $A_d A_e \subseteq A_{d + e}$ for all $d, e\geq 0$.
		
		\item If $R$ is any ring, a graded ring $A$ is called a \inlinedef{graded $R$-algebra} if $A_0$ is an $R$-algebra.
		
		\item A \inlinedef{homomorphism $f \colon A \to B$ of graded $R$-algebras} $A, B$ is an $R$-algebra homomorphism such that $f(A_d) \subseteq B_d$ for all $d \geq 0$.
	\end{enumerate}
\end{definition}

\begin{remark}
	Let $A$ be a graded ring.
	\begin{enumerate}
		\item The ring $A$ is an $A_0$-algebra. Indeed, since $A_0A_0 \subseteq A_0$ it follows that $A_0$ is a ring.
		
		\item Let $R$ be any ring. Then $A$ is a graded $R$-algebra if and only if $A$ is an $R$-algebra and $RA_d = A_d$ for all $d$.
	\end{enumerate}
\end{remark}

\begin{example}
	A prototypical example of a graded ring is the polynomial ring $A = k[x]$ for a field $k$. The grading on $A$ is the grading by degrees, that is,
	\begin{equation*}
		A_d \coloneqq \Set{ax^d : a \in k}
	\end{equation*}
	for $d \geq 0$. Note that $A_0 = k$.
\end{example}

If $A$ is a graded $R$-algebra and $\delta \geq 0$ is an integer, we can define a new $R$-algebra from $A$ by
\begin{equation*}
	A^{(\delta)} \coloneqq \oplus_{d \geq 0} A_{d \delta}.
\end{equation*}
This is a graded $R$-algebra. We now turn to graded modules.

\begin{definition}
	Let $A$ be a graded ring. A \inlinedef{graded $A$-module} is an $A$-module $M$ with a decomposition $M = \oplus_{d \in \NN} M_d$ of Abelian groups such that $A_d M_e \subseteq M_{d + e}$ for all $d \geq 0$ and $e \in \ZZ$.
\end{definition}

Elements of $A_d$ and $M_d$ are called \inlinedef{homogeneous of degree $d$}. If $a \neq 0$ is homogeneous of degree $d$, then its degree $\deg a$ is defined to be $d$. It is then evident that every element $x \in A$ can be uniquely written in the form
\begin{equation*}
	x = \sum_{d} x_d
\end{equation*}
where $x_d \in A_d$ and $x_d = 0$ for almost all $d$. A similar statement holds for elements of a graded $A$-module. The terms $x_d$ are called \inlinedef{homogeneous components} of $x$.

We now define homomorphisms of graded $A$-modules. A homomorphism $f \colon M \to N$ of $A$-modules is \inlinedef{graded} if $f(M_d) \subseteq N_d$ for all $d$. More generally, a homomorphism of $A$-modules $u \colon M \to N$ is \inlinedef{homogeneous of degree $k$} if $u(M_{d}) \subseteq N_{d + k}$ for all $d$.

\subsection*{Homogeneous Submodules and Homogeneous Ideals}

\begin{definition}
	Let $A$ be a graded ring and let $M$ be a graded $A$-module. A submodule $N \subseteq M$ is said to be \inlinedef{homogeneous} if it satisfies one of the following equivalent conditions:
	\begin{enumerate}
		\item $N = \oplus_{d} (N \cap M_d)$.
		\item $N$ is generated by homogeneous elements of $M$.
		\item For every $x \in N$, the homogeneous components of $x$ are in $N$.
	\end{enumerate}
	
	A \inlinedef{homogeneous ideal} of $A$ is an ideal which is also a homogeneous submodule of $A$ viewed as a graded $A$-module.
\end{definition}

We show that the conditions are indeed equivalent.
\begin{proof}
	(1.) $\Rightarrow$ (2.) Take the generators of $N$ to be elements of $N \cap M_d$.
	
	(2.) $\Rightarrow$ (3.) Suppose $x \in M$. Then by assumption $x = \sum_{i = 1}^N n_i$ for suitable $N \in \NN$ and suitable homogeneous elements $n_i \in N$. Since the sum of homogeneous elements of the same degree $d$ is still homogeneous of degree $d$, we may collect those terms and assume that $n_i$ is homogeneous of degree $i$. But such expression is unique, and so the homogeneous components of $x$ are indeed the $n_i$, which are elements of $N$.
	
	(3.) $\Rightarrow$ (1.) Suppose that $x \in N$. Write $x = \sum_d x_d$ in terms of its homogeneous components $x_d \in M_d$. By assumption, $x_d \in N$. Hence $x_d \in N \cap M_d$. Uniqueness follows immediately. This proves that $N = \oplus_{d} (N \cap M_d)$.
\end{proof}

From (1.) it is clear that $N$ is also a graded $A$-module.	The quotient of a graded module by its homogeneous submodule inherits the grading, as well as the sum and the intersection of a family of graded submodules.

\begin{proposition}
	Let $M$ be a graded $A$-module, $N$ a graded submodule of $M$, and let $(N_i)_{i \in I}$ be a family of graded $A$-submodules of $M$. 
	\begin{enumerate}
		\item The quotient module $M/N$ is graded via the grading $(M/N)_d = M_d/(N \cap M_d)$.
		
		\item The sum $\sum_{i \in I} N_i$ and the intersection $\bigcap_{i \in I} N_i$ are graded $A$-submodules of $M$.
	\end{enumerate}
\end{proposition}
\begin{proof}
	\todo{TODO.}
\end{proof}

\begin{definition}
	Let $M$ be a graded $A$-module and $S \subseteq M$ a subset. The \inlinedef{homogeneous submodule generated by $S$} is defined to be the intersection of all homogeneous submodule $N \subseteq M$ such that $S \subseteq N$.
\end{definition}

The above proposition guarantees that the resulting intersection is indeed homogeneous. Furthermore, a submodule $N$ is homogeneous if and only if the homogeneous submodule generated by $N$ is $N$ itself. If $I \subseteq A$ is an ideal of $A$, then we denote by $I^h$ the homogeneous ideal generated by $I$. 

\subsection*{Graded Localisation}

We discuss a localisation of 

\subsection*{Relevant Prime Ideals}

For graded rings, only a special type of ideals plays a role.

\begin{remark}
	Let $A$ be a graded ring. The ideal $A_+ \coloneqq \oplus_{d \geq 1} A_d$ is a homogeneous ideal because it is generated by homogeneous elements.
\end{remark}

The ideal $A_+$ is called the \inlinedef{irrelevant ideal}. More generally, for any ideal $I$ of $A$, we define
\begin{equation*}
	I_+ \coloneqq I \cap A_+.
\end{equation*} 
A homogeneous prime ideal $\mathfrak{p}$ of $A$ is \inlinedef{relevant} if it does not contain the irrelevant ideal. This is equivalent to $\mathfrak{p}_+ \subsetneq A_+$. Indeed, if $A_+ \subseteq \mathfrak{p}$, then $\mathfrak{p}_+ = \mathfrak{p} \cap A_+ = A_+$. Conversely, if $A_+ \not\subseteq \mathfrak{p}$, then there is $x \in A_+ \setminus \mathfrak{p}$, and so $x \notin \mathfrak{p}_+$.



\section{Projective Spectrum}

\begin{definition}
	Let $A$ be a graded ring. The set of all relevant prime ideals of $A$ is called the \inlinedef{projective spectrum} of $A$. It is denoted by $\Proj A$.
\end{definition}

We now endow a topology on $\Proj A$ in a similar manner to the Zariski topology. Define, for each homogeneous ideal $I \subseteq A$,
\begin{equation*}
	V_+(I) \coloneqq \Set*{\mathfrak{p} \in \Proj A : \mathfrak{p} \subseteq I}.
\end{equation*}
The following are as easy to prove as for spectra of a ring.
\begin{enumerate}
	\item $V_+(0) = \Proj A$ and $V_+(1) = \emptyset$.
	\item $\bigcap_i V_+(I_i) = V_+(\sum_i I_i)$
	\item $V_+(I) \cup V_+(J) = V_+(I \cap J)$.
\end{enumerate}
Therefore the sets of the form $V_+(I)$, where $I$ runs through all homogeneous ideals of $A$, form the closed sets of a topology on $\Proj A$. If $f \in A$, the projective principal open set of $f$ is the open set 
\begin{equation*}
	D_+(f) = (\Proj A) \setminus V_+(f).
\end{equation*}

\begin{proposition}
	The projective spectrum $\Proj A$ of a graded ring $A$ is a subspace of $\Spec A$.
\end{proposition}
\begin{proof}
	Clearly $\Proj A \subseteq \Spec A$ as sets. We note that if $\mathfrak{a}$ is an ideal of $A$, then
	\begin{equation*}
		V(\mathfrak{a}) \cap \Proj A = V_+(\mathfrak{a}^h),
	\end{equation*}
	where $\mathfrak{a}^h$ is the homogeneous ideal generated by $\mathfrak{a}$. Thus the topology on $\Proj A$ is the subspace topology.
\end{proof}


